% Problem dossier: VI.06.P8
\subsection*{Problem VI.06.P8 --- Radical versus prime: one locus or several?}
\label{prob:vi-06-p8}

\paragraph{Problem.}
Let \(k\) be a field and
\[
I=(xy)\subset k[x,y].
\]
\begin{enumerate}
\item Prove that \(I=(x)\cap(y)\).
\item Deduce that \(I\) is radical.
\item Prove that \(I\) is not prime.
\item Explain geometrically why radicality corresponds to ``no nilpotent thickening'' while
primality corresponds to ``one irreducible component.''
\end{enumerate}

\ifdefined\FullProblemDossiers
\paragraph{Why this problem matters.}
Prime and radical ideals are often confused because every prime is radical.  The crossing
axes give the cleanest counterexample to the converse.

\paragraph{Useful ideas before starting.}
Divisibility in the UFD \(k[x,y]\), intersections of radical ideals, and the quotient-domain
criterion.

\paragraph{Guided hints.}
If a polynomial lies in both \((x)\) and \((y)\), use coprimeness of \(x\) and \(y\) to
show divisibility by \(xy\).

\paragraph{Solution.}
The inclusion \((xy)\subseteq(x)\cap(y)\) is immediate.  Conversely, if \(f\) is divisible
by both \(x\) and \(y\), coprimeness of the irreducibles \(x,y\) implies that \(xy\mid f\).
Thus
\[
(xy)=(x)\cap(y).
\]
Both \((x)\) and \((y)\) are prime and hence radical; their intersection is radical.
However, \(xy\in I\) while \(x,y\notin I\), so \(I\) is not prime.

Geometrically, \(V(I)\) is the reduced union of two axes.  Radicality says the equations do
not carry an infinitesimal thickening; nonprimality says the locus splits into more than one
irreducible piece.

\paragraph{What happened mathematically.}
The chain
\[
\text{maximal}\Rightarrow\text{prime}\Rightarrow\text{radical}
\]
reflects increasingly weaker geometric conditions: closed point, irreducible locus, reduced
closed set.

\paragraph{Extensions.}
For pairwise nonassociate irreducible polynomials \(f_1,\dots,f_r\), investigate when
\((f_1\cdots f_r)\) is radical.

\paragraph{Provenance.}
New canonical synthesis from the chapter's prime/radical comparison.
\fi
